\section*{الفصل \ref{ch:g11:igneous-rocks} --- الصخور النارية}

\begin{solution}{exo:g11:igneous-rocks:1}
\omterm{def:g11:igneous-rocks:volcanic-plutonic}{الصخر البركاني} تكوّن من صهارة تصلبت على السطح، لابةً؛ \omterm{def:g11:igneous-rocks:volcanic-plutonic}{والصخر
الجوفي} تبلور في العمق، داخل صخور أخرى. والأزواج: البازلت
\omterm{def:g11:ocean-ridges:gabbro}{والغابرو} (\omterm{def:g11:igneous-rocks:felsic-mafic}{المافيان})، \omterm{def:g11:igneous-rocks:rock-names}{والأنديزيت} \omterm{def:g11:igneous-rocks:rock-names}{والديوريت} (المتوسطان)، \omterm{def:g11:igneous-rocks:rock-names}{والريوليت} والغرانيت
(\omterm{def:g11:igneous-rocks:felsic-mafic}{الفلسيان}).
\end{solution}

\begin{solution}{exo:g11:igneous-rocks:2}
(1) زجاجي. (2) خشن الحبيبات. (3) \omterm{def:g11:igneous-rocks:porphyritic}{بورفيري} (بلورات أوليفين ظاهرة في
أرضية دقيقة). (4) دقيق الحبيبات.
\end{solution}

\begin{solution}{exo:g11:igneous-rocks:3}
\qty{48}{\percent}: \omterm{def:g11:igneous-rocks:felsic-mafic}{مافية}، \omterm{def:g11:ocean-ridges:gabbro}{غابرو}. \qty{60}{\percent}: \omterm{def:g11:igneous-rocks:felsic-mafic}{متوسطة}، \omterm{def:g11:igneous-rocks:rock-names}{ديوريت}.
\qty{72}{\percent}: \omterm{def:g11:igneous-rocks:felsic-mafic}{فلسية}، غرانيت.
\end{solution}

\begin{solution}{exo:g11:igneous-rocks:4}
\omterm{def:g11:igneous-rocks:intrusion}{القاطع} صفيحة من الصخر الناري تقطع طبقات الصخور التي
تداخلت فيها؛ \omterm{def:g11:igneous-rocks:intrusion}{والعتبة} صفيحة موازية للطبقات، واقعة بينها؛ \omterm{def:g11:igneous-rocks:intrusion}{والبلوتون}
\omterm{def:g11:igneous-rocks:intrusion}{جسم متداخل} كبير أيًا كان شكله تكوّن في العمق. \omterm{def:g11:igneous-rocks:intrusion}{والقاطع} الرأسي الذي يصعد
من \omterm{def:g11:igneous-rocks:intrusion}{البلوتون} هو الذي غذّى البركان.
\end{solution}

\begin{solution}{exo:g11:igneous-rocks:5}
يتكوّن الأوليفين أولًا؛ وفلسبار البوتاسيوم والكوارتز وأحيانًا المسكوفيت
أخيرًا. ويبدأ البلاجيوكليز غنيًا بالكالسيوم ثم يصير أغنى فأغنى
بالصوديوم كلما بردت الصهارة.
\end{solution}

\begin{solution}{exo:g11:igneous-rocks:6}
(1) A: \qty{74}{\percent} من \ce{SiO2}، \omterm{def:g11:igneous-rocks:felsic-mafic}{فلسي}؛ وB: \qty{47}{\percent}، \omterm{def:g11:igneous-rocks:felsic-mafic}{مافي}.
(2) A فقير بالأكاسيد \ce{FeO} و\ce{MgO} و\ce{CaO} وغني بالأكسيدين \ce{Na2O}
و\ce{K2O}، كما هي الصهارة \omterm{def:g11:igneous-rocks:felsic-mafic}{الفلسية}؛ وB غني بالأكاسيد \ce{FeO} و\ce{MgO} و\ce{CaO}،
فقير بالأكسيد \ce{Na2O}، كما هي الصهارة \omterm{def:g11:igneous-rocks:felsic-mafic}{المافية}. (3) A \omterm{def:g11:igneous-rocks:rock-names}{ريوليت}؛ وB \omterm{def:g11:ocean-ridges:gabbro}{غابرو}.
\end{solution}

\begin{solution}{exo:g11:igneous-rocks:7}
يتناسب الزمن مع مربع السمك. \qty{1}{m}: $10/4 = 2.5$
يوم. \qty{20}{m}: $10 \times 10^2 = 1000$ يوم، أي نحو 2.7 سنة.
\qty{200}{m}: $10 \times 100^2 = \num{100000}$ يوم، أي نحو 270 سنة.
\end{solution}

\begin{solution}{exo:g11:igneous-rocks:8}
(1) قسمة \unit{m^2} على \unit{m^2/s} تعطي ثواني. (2) $400/10^{-6} =
\qty{4e8}{s}$، أي $4 \times 10^8 / 3.16 \times 10^7 \approx 12.7$
سنة. (3) $0.21 \times 12.7 \approx 2.7$ سنة.
\end{solution}

\begin{solution}{exo:g11:igneous-rocks:9}
إنه \omterm{def:g11:igneous-rocks:porphyritic}{بورفيري}: بلورات بلاجيوكليز ظاهرة في أرضية دقيقة الحبيبات.
فقد بردت الصهارة أولًا ببطء في العمق، ونمت بلورات البلاجيوكليز
كبيرة بينما بقي معظم الصهارة سائلًا؛ ثم صعدت وبردت بسرعة،
فتبلور السائل المتبقي أرضيةً دقيقة.
\end{solution}

\begin{solution}{exo:g11:igneous-rocks:10}
(1) في الطرف الأيسر: 2 Mg و2 Si و6 O؛ وفي الطرف الأيمن: 2 Mg و2 Si و6 O. (2) \qty{1}{kg} من
الأوليفين يساوي $1000/140.7 \approx 7.11$~\unit{mol}، فيأخذ العدد نفسه من مولات
السيليكا: $7.11 \times 60.1 \approx \qty{427}{g}$. (3) إذا لم تبق سيليكا في
السائل، فلا يمكن للتفاعل أن يستمر: فيبقى الأوليفين.
\end{solution}

\begin{solution}{exo:g11:igneous-rocks:11}
(1) $2.5 \times 20 \times 1100 = \qty{55000}{MJ}$، أي \qty{5.5e10}{J}.
(2) $900 \times 20 = \qty{18000}{MJ}$، أي \qty{1.8e10}{J}. (3) حرارة
التبلور نحو ثلث الحرارة المفقودة بالتبرد، وربع
المجموع: وهذا يكفي لإبطاء التجمد كثيرًا، لأنها تنطلق كلها
فوق خط التصلب.
\end{solution}

\begin{solution}{exo:g11:igneous-rocks:12}
(1) أصغر البلورات عند التماس، الذي تجمد في أقل من يوم؛
وأكبرها في المركز، الذي بقي سائلًا نحو 950 يومًا. (2) لا: فأول
\qty{5}{m} استغرق نحو 545 يومًا، وآخر \qty{5}{m} نحو 400 يوم؛ فالجبهة
تتباطأ في البداية، ثم تتسارع قرب المركز، حيث تلتقي الحرارة
الخارجة من الجانبين. (3) حتى في مركزها تجمدت \omterm{def:g11:igneous-rocks:intrusion}{العتبة} في
أقل من ثلاث سنوات؛ أما الغرانيت ذو البلورات الكبيرة فيحتاج إلى زمن أطول بكثير.
\end{solution}

\begin{solution}{exo:g11:igneous-rocks:13}
(1) \omterm{def:g11:igneous-rocks:intrusion}{البلوتون} أسمك بمئة مرة: $2.6 \times 100^2 = \num{26000}$
سنة. (2) على سبيل المثال: \omterm{def:g11:igneous-rocks:intrusion}{البلوتون} العميق يقع في صخر ساخن أصلًا، فهو
يفقد حرارته أبطأ مما لو كان في صخر عند \qty{50}{\degreeCelsius}؛ وقد
يتغذى بصهارة جديدة مرة بعد مرة؛ \omterm{def:g11:igneous-rocks:intrusion}{والبلوتون} ليس صفيحة لا نهائية؛
وخواص الصخر تتغير مع درجة الحرارة.
\end{solution}

\begin{solution}{exo:g11:igneous-rocks:14}
الصهارة \omterm{def:g11:igneous-rocks:felsic-mafic}{المافية} ساخنة وفقيرة بالسيليكا: فهي تبلور المعادن الأولى
من السلسلة --- الأوليفين والبيروكسين والبلاجيوكليز الغني بالكالسيوم --- وتنفد منها
السيليكا قبل أن يتكوّن الكوارتز. أما الصهارة \omterm{def:g11:igneous-rocks:felsic-mafic}{الفلسية} فأبرد وغنية
بالسيليكا: لا يكون الأوليفين مستقرًا فيها، وهي تمضي حتى نهاية
السلسلة، فتكوّن البلاجيوكليز الغني بالصوديوم وفلسبار البوتاسيوم والكوارتز.
\end{solution}

\begin{solution}{exo:g11:igneous-rocks:15}
(1) $b = \sqrt{t\kappa/C} = \sqrt{86400 \times 10^{-6}/0.21} \approx
\qty{0.64}{m}$. (2) يعطي القانون $0.21 \times 25/10^{-6} \approx
\qty{5.3e6}{s}$، أي نحو شهرين: من رتبة الشهر نفسها (فالمسيل اللابي
يفقد الحرارة أيضًا في الهواء من أعلاه). (3) زجاجي للمواد المقذوفة في
الهواء أو الماء؛ ودقيق الحبيبات للمسيلات اللابية والصفائح الرقيقة؛ وأخشن، لكنه يبقى
صغير الحبيبات، \omterm{def:g11:igneous-rocks:intrusion}{للأجسام المتداخلة} الضحلة السميكة.
\end{solution}

\begin{solution}{pb:g11:igneous-rocks:1}
\textbf{1.} \omterm{def:g11:igneous-rocks:intrusion}{عتبة}: فهي موازية للطبقات، واقعة بينها.

\textbf{2.} عند التماس، دقيق الحبيبات جدًا، يكاد يكون زجاجيًا؛ وفي المركز،
دقيق إلى متوسط الحبيبات، تُرى المعادن بالكاد بالعدسة.

\textbf{3.} فقد التماس حرارته بأسرع ما يكون، في الحجر الرملي البارد: فلم يكن
للبلورات وقت لتنمو.

\textbf{4.} إنها حافة مبردة. وهي تثبت أن الصهارة كانت ساخنة حين استقرت
في موضعها، وأن الحجر الرملي كان باردًا.

\textbf{5.} \omterm{def:g11:igneous-rocks:intrusion}{قاطع}، هو مغذّي \omterm{def:g11:igneous-rocks:intrusion}{العتبة}.

\textbf{6.} \omterm{def:g11:igneous-rocks:felsic-mafic}{مافية}: \qty{49}{\percent} من \ce{SiO2}، وغنية بالحديد والمغنيسيوم
والكالسيوم.

\textbf{7.} الأوليفين والبيروكسين والبلاجيوكليز الغني بالكالسيوم.

\textbf{8.} بازلت؛ \omterm{def:g11:igneous-rocks:volcanic-plutonic}{والصخر الجوفي} ذو التركيب نفسه هو \omterm{def:g11:ocean-ridges:gabbro}{الغابرو}.

\textbf{9.} يتوقف \omterm{def:g11:igneous-rocks:texture}{النسيج} على معدل التبرد، الذي يختلف من
التماس إلى المركز، لا على التركيب.

\textbf{10.} الأوليفين أولًا؛ ثم يبدأ البلاجيوكليز تقريبًا عند النقطة نفسها التي يبدأ عندها
البيروكسين.

\textbf{11.} $b^2/\kappa = 400/10^{-6} = \qty{4e8}{s}$.

\textbf{12.} $0.206 \times 4 \times 10^8 \approx \qty{8.2e7}{s}$، أي
$8.2 \times 10^7 / 3.16 \times 10^7 \approx 2.6$ سنة.

\textbf{13.} $0.206 \times 4/10^{-6} \approx \qty{8.2e5}{s}$، أي نحو 9.5
يوم.

\textbf{14.} $0.07 \times 4 \times 10^8 = \qty{2.8e7}{s}$، أي نحو 0.9 سنة.
فالصهارة المتبلورة تطلق حرارة، يجب أن تُصرَّف هي أيضًا قبل
أن تصير الصهارة صلبة.

\textbf{15.} نحو 950 يومًا مقابل 0.6 يوم: أي أسرع بنحو 1600 مرة.

\textbf{16.} نحو 950 يومًا، أي نحو 2.6 سنة: وهذا أطول بكثير من المسيل اللابي،
وضمن الأشهر إلى السنوات التي \omterm{def:g11:igneous-rocks:intrusion}{للأجسام المتداخلة} الضحلة.

\textbf{17.} الزمن الأطول ألف مرة يحتاج إلى سمك أكبر بنحو $\sqrt{1000} \approx
32$ مرة: أي نحو \qty{630}{m}.

\textbf{18.} قرب السطح، كان جسم من الصهارة سينبثق، أو سيبرد بسرعة كبيرة؛
أما الجسم الكبير الذي يبرد زمنًا طويلًا جدًا فلا بد أن يكون عميقًا، في صخر
ساخن أصلًا.

\textbf{19.} لقد أزالتها التعرية.

\textbf{20.} استغرقت \omterm{def:g11:igneous-rocks:intrusion}{العتبة} نحو \textbf{2.6 سنة} لتتصلب.
\end{solution}
